\documentclass{ximera}
%geschreven 12/11/26

\title{Natuurconstanten en gegevens}
\author{Daan Lipkens}

\begin{document}
\begin{abstract}
    Een overzicht van de constanten en deeltjesgegevens die je in deze cursus nodig hebt.
\end{abstract}
\maketitle
\label{hoofdstuk:constanten}

\section{Natuurconstanten}\label{sec:natuurconstanten}
In de oefeningen gebruik je telkens dezelfde waarden voor de natuurconstanten. Je vindt ze samen in tabel~\ref{tab:natuurconstanten}. De waarden staan op drie beduidende cijfers, wat voldoende is voor alle oefeningen in deze cursus.

\begin{table}[ht!]
    \begin{center}
        \begin{tabular}{l c l}
            \toprule
            Grootheid & Symbool & Waarde \\
            \midrule
            constante van Coulomb (vacuüm of lucht) & $k$ & $8{,}99 \cdot 10^{9} \, \mathrm{N \cdot m^2/C^2}$ \\
            elementaire lading & $e$ & $1{,}60 \cdot 10^{-19} \, \mathrm{C}$ \\
            magnetische permeabiliteit van vacuüm & $\mu_0$ & $4\pi \cdot 10^{-7} \, \mathrm{T \cdot m/A} = 1{,}26 \cdot 10^{-6} \, \mathrm{T \cdot m/A}$ \\
            zwaarteveldsterkte aan het aardoppervlak & $g$ & $9{,}81 \, \mathrm{N/kg}$ \\
            \bottomrule
        \end{tabular}
    \end{center}
    \caption{Natuurconstanten die in deze cursus gebruikt worden.}
    \label{tab:natuurconstanten}
\end{table}

\begin{remark}[title={Elementaire lading en teken}]\nl
    De elementaire lading $e$ is altijd positief. Een proton heeft een lading $+e$, een elektron een lading $-e = -1{,}60 \cdot 10^{-19} \, \mathrm{C}$. Het teken zet je dus zelf bij de lading van het deeltje.
\end{remark}

\begin{remark}[title={Zwaarteveldsterkte en valversnelling}]\nl
    De eenheid $\mathrm{N/kg}$ is gelijk aan $\mathrm{m/s^2}$. Je vindt $g$ daarom ook terug als de valversnelling $g = 9{,}81 \, \mathrm{m/s^2}$.
\end{remark}

\section{Eigenschappen van deeltjes}\label{sec:deeltjes}
Tabel~\ref{tab:deeltjes} geeft de lading en de massa van de deeltjes waaruit een atoom bestaat. Dezelfde gegevens vind je ook in het hoofdstuk over elektrische krachtwerking, zie tabel~\verwijs[in Ladingen]{tabel:ladingen_proton_elektron}{elektrische_krachtwerking_lading}.

\begin{table}[ht!]
    \begin{center}
        \begin{tabular}{l c r r}
            \toprule
            Deeltje & Lading & Lading ($\mathrm{C}$) & Massa ($\mathrm{kg}$) \\
            \midrule
            proton & $+e$ & $+1{,}60 \cdot 10^{-19}$ & $1{,}673 \cdot 10^{-27}$ \\
            neutron & $0$ & $0$ & $1{,}675 \cdot 10^{-27}$ \\
            elektron & $-e$ & $-1{,}60 \cdot 10^{-19}$ & $9{,}109 \cdot 10^{-31}$ \\
            \bottomrule
        \end{tabular}
    \end{center}
    \caption{Lading en massa van proton, neutron en elektron.}
    \label{tab:deeltjes}
\end{table}

\begin{remark}[title={Wat moet je kennen?}]\nl
    De lading van proton en elektron ken je vanbuiten. De massa's hoef je niet te kennen: die worden gegeven of je zoekt ze op in deze tabel. In oefeningen rond je de massa's af op drie beduidende cijfers, bv. $m_p = 1{,}67 \cdot 10^{-27} \, \mathrm{kg}$ en $m_e = 9{,}11 \cdot 10^{-31} \, \mathrm{kg}$.
\end{remark}

\section{Materiaalgegevens}\label{sec:materiaalgegevens}
Sommige gegevens hangen af van de stof. De relatieve permeabiliteit $\mu_r$ van enkele stoffen vind je in tabel~\verwijs[in Elektromagneten]{tab:mu_r}{magnetisme_elektromagneet}.

\end{document}
